\documentclass[10pt]{article}
\usepackage[margin=20mm]{geometry}
\usepackage{amsmath,amssymb,fontspec,unicode-math,xurl,hyperref}
\IfFontExistsTF{DejaVu Serif}{\setmainfont{DejaVu Serif}}{\setmainfont{Latin Modern Roman}}
\setmathfont{latinmodern-math.otf}
\setlength{\parindent}{0pt}\setlength{\parskip}{6pt}\sloppy
\begin{document}
Abstract: Independent restricted-dual-norm and transport calculations support the main regression and classification reformulations. The interpretation as an unconditional out-of-sample guarantee is false; several proof formulas and empirical comparisons need correction. These are assessed as Minor because the central finite-dimensional equivalences survive direct repair, and the paper explicitly leaves statistical guarantees to future work. No author implementation or experimental reproduction was executed.\par
Author: AI-assisted technical review\par
Date:\par
Title: Technical Verification of Knowledge-Guided Wasserstein Distributionally Robust Optimization\par
\# Identity and scope\par
Zitao Wang, Ziyuan Wang, Molei Liu and Nian Si, ICML 2025, I25-43697 (0328), Knowledge-Guided Wasserstein Distributionally Robust Optimization. The frozen publisher PDF contains 23 pages. All main scientific text, pp. 1–9, and Appendices A–F, pp. 11–23, were read. The bibliography was not independently audited. Both original figures and Tables 2–5 were visually checked, along with decisive formulas. No relevant repository was identified in the bounded source acquisition; separate OpenReview attachments were inaccessible. Later errata and raw experimental results were unavailable to this review. Page numbers below are one-based PDF numbers, equal to printed numbers.\par
\# Core contribution and independent checks\par
Equation reference: eq:dual. Strong transfer constrains perturbations \(\Delta\) to \(S=\{\Delta:\theta_m^T\Delta=0\ \forall m\}=\Theta^\perp\), where \(\Theta\) is the span of the prior predictors. The restricted support function is \[
 a(\beta)=\sup_{u\in S,\,\|u\|_q\leq1}|\beta^Tu|
 =\inf_{\vartheta\in\Theta}\|\beta-\vartheta\|_p.\] This is the dual norm on the quotient by \(\Theta\), including \(p=1,\infty\). For squared loss and average transport cost at most \(\delta\), Minkowski and [eq:dual] (\#eq:dual) bound root risk by \(\sqrt{\mathrm{MSE}_N}+\sqrt\delta\,a(\beta)\). For nonzero empirical residual vector \(r\), choose an attaining unit \(u\in S\) with \(\beta^Tu=a\) and \(\Delta_i=-\sqrt\delta\,r_i u/\sqrt{\mathrm{MSE}_N}\). This attains the bound and uses exactly budget \(\delta\). When every residual is zero, a shift of length \(\sqrt\delta\) attains \(\delta a^2\); \(a=0\) is immediate. Thus the fixed-\(\beta\) identity, and hence Theorem 3.3, hold without the appendix’s singular divisions.\par
For logistic and hinge losses, the Lipschitz upper bound gives empirical loss plus \(\delta a\). A mass of order \(\delta/t\) sent distance \(t\) along an adverse unit direction approaches that bound as \(t\to\infty\), because both losses have asymptotic adverse slope one. This verifies the supremum formula behind Theorem 3.6, including the endpoint norms. It need not produce an attained worst-case distribution.\par
For weak transfer, \(A=I+\lambda\theta\theta^T\) and the correct penalty is \[\sqrt{\beta^TA^{-1}\beta},\qquad
 A^{-1}=I-\frac{\theta\theta^T}{\|\theta\|_2^2+\lambda^{-1}}.\] This equals \(\|\beta\|_{\Psi_{\lambda^{-1}}}\) under the main-text definition of \(\Psi\). It establishes Theorem 3.5 by the same argument. Replacing \(I\) by positive definite \(\Lambda\) verifies Corollaries 3.7 and D.1. The full-rank prior span reduces to ERM, as the authors explicitly note. Independent numerical checks use \(\theta=(1,2)\), \(\beta=(3,-1)\): \(a=3.5,7/\sqrt5,7/3\) for \(p=1,2,\infty\). A three-residual example also checks exact transport budget and attained risk. These are checks of identities, not reruns of reported learning experiments.\par
\# Localized findings and repairs\par
Problem 1 (The sample-outside-the-ambiguity-set interpretation is unsupported).  The sentence following Proposition 3.1, p. 5, describes its minimax bound as out-of-sample performance at least as good as the empirical domain adapter. This is false without a condition placing the target law in the ambiguity set. Severity: Minor, a necessary scope correction.\par
Let \(d=2\), \(\theta=e_1\), \(N=1\), and observe \((x,y)=(e_1,1)\). For any \(\delta>0\), the root robust objective is \(|1-\beta_1|+\sqrt\delta|\beta_2|\), with unique minimizer \(\beta=e_1\) and value zero. Let the true target distribution give equal probability to \((e_1,1)\) and \((e_2,1)\). The observed sample has probability \(1/2\), but this predictor has target squared risk \(1/2\), while the empirical adapter risk is zero. The target law cannot belong to the ambiguity set: its \(\theta^TX\) marginal differs from the frozen empirical marginal.\par
Proposition 3.1 itself is valid: choose \(\beta=\alpha\theta\) and use the preserved joint law of \((Y,\theta^TX)\). No minimax interchange is needed. Its guarantee is for laws in that restricted set. For a continuous \((Y,\theta^TX)\) population, the exact preservation of the empirical discrete marginal generally precludes population membership for every radius. Increasing \(\delta\) does not cure that obstruction. The text should state this scope, distinguish the limit \(\beta\in\Theta\) from \(\beta\to\theta\), and avoid a general no-negative-transfer interpretation. The conclusion’s explicit deferral of statistical guarantees is important counterevidence to treating this as a failed central statistical theorem.\par
Problem 2 (Several displayed proof formulas require local corrections).  Appendix B, pp. 17–18, Appendix C, pp. 18–19, and Proposition E.2(4), p. 22, contain algebraic errors. Severity: Minor; the main equivalences have independent derivations above.\par
In the weak-transfer proof, write \(r=y-\beta^Tx\) and \(b=\sqrt{\beta^TA^{-1}\beta}\). The correct radial maximization is \[\sup_{t\geq0}\{(b^2-\gamma)t^2+2|r|bt\}
 =\frac{r^2b^2}{\gamma-b^2}\quad(\gamma>b^2).\] The printed negative linear term is wrong, the next line drops the square on \(b\), and the finiteness threshold must compare \(\gamma\) with \(b^2\). The residual is also inadvertently defined as squared on p. 17. For \(\theta=e_1\), \(\beta=e_2\), \(r=1\), \(\gamma=2\), the actual inner supremum is \(1\), attained at \(\Delta=-e_2\), whereas the printed radial expression \(-t^2-2t\) has supremum zero. For \(b>0\), at \(\gamma=b^2\) a nonzero residual gives infinity; zero-residual and \(b=0\) cases require separate limits, also in the proof of Theorem 3.3.\par
For \(h(x)=\log(1+e^{-x})\), the correct conjugate has domain \([-1,0]\): \[h^*(z)=(-z)\log(-z)+(1+z)\log(1+z),\quad -1\leq z\leq0.\] The printed \([0,1]\) formula belongs to \(\log(1+e^x)\). At \(z=1/2\), \(zx-h(x)\to\infty\) as \(x\to\infty\), contradicting the finite printed value. Lemma C.1 must use \(-\alpha\beta\) in its biconjugate expression, or redefine \(h\) consistently. The final distance to the linear span is unchanged by this sign, so its conclusion survives.\par
The signed linear “weighted penalty” in Section 3.4, p. 7, must instead be the square root of \(\min_\kappa\{(\beta_1-\kappa\theta_1)^2/d_1+
(\beta_2-\kappa\theta_2)^2/d_2\}\). As printed the linear minimization can equal \(-\infty\). Theorem E.5’s unrestricted sum-conjugacy statement generally needs closure or a qualification; the finite continuous norm functions used here supply the needed qualification. The rank-one difference between the weak metric and the strong projection is positive semidefinite, not positive definite: equality holds for \(\beta\perp\theta\), contrary to the strict comparison on p. 6. The ellipse and the non-strict inclusion remain correct.\par
Problem 3 (Empirical prose and plotting do not consistently match the tables).  Section 4.1, Figure 2 and Tables 3–5 need reconciliation. Severity: Minor; the reported improvements in several regimes remain supported.\par
At \(N=80,s=1,\rho=0.3\), Table 3 gives KG-WDRO accuracy \(0.583\) versus WDRO \(0.654\), a loss of \(7.1\) percentage points. At \(N=50,s=1\) and the same correlation it gives \(0.586\) versus \(0.619\). Thus the unqualified claim that both transfer methods significantly outperform target-only WDRO is contradicted by the paper’s own averages. Tables give no dispersion or paired confidence intervals, so statistical significance cannot be inferred from their rounded means. This does not negate the clear favorable regimes at high correlation or the low-sample comparisons.\par
The two lower panels of Figure 2 label \(N=70\) but draw the WDRO dashed baseline near \(0.49\), whereas Table 5 gives \(0.375\) and \(0.354\) for the displayed settings. The latter values are below the plotted vertical minimum \(0.4\). This is a visible inconsistency, not a claim about which underlying experiment is correct. The generating scheme is described clearly enough to understand the comparisons, but Appendix A.2.3 lists a correlation grid beginning at \(0.3\), while Section 4.3 and Table 5 start at \(0.7\) and include \(0.75,0.85\); the actual grid needs identification. By contrast, the real-data table does support five wins in eight states and a \(7.587\%\) reduction in aggregate log loss, rounding to \(7.6\%\). Its evaluation split and uncertainty remain insufficiently specified for independent reproduction without code/data-processing details.\par
\# Assessment, evidence and limits\par
The central transport-to-regularization mechanism is sound under its stated finite-sample ambiguity sets. Its handling of collinearity permits scalar rescaling and sign changes, but does not by itself guarantee invariance to arbitrary feature-specific transformations or eliminate negative transfer. The mathematical repairs above are local and do not establish a new population generalization theorem. Accordingly the overall assessment is Minor, with no demonstrated Major defect in the principal reformulations. All judgments remain pending; author adoption is unknown and no human review is recorded. Focused independent acceptance is documented in the existing verification files.\par
1 Z. Wang, Z. Wang, M. Liu and N. Si, Knowledge-Guided Wasserstein Distributionally Robust Optimization, ICML 2025, frozen publisher PDF. https://proceedings.mlr.press/v267/wang25df.html.\par
\# Reproducibility\par
Source SHA-256: `4ee170ac69376b3133a4b4a790342b5c194c249be02af44e579aba38974ca10b`. The independent Python script uses NumPy and SciPy for the low-dimensional dual-norm, inverse-matrix, transport, conjugate and table checks. Inputs, outputs, exact source-page hashes and initial visual observations accompany this report. No author modules were imported or executed. The final TeX is actually compiled and its PDF pages inspected, with build and page-image records saved separately. Selected reviewer: GPT-6 Astra / Max; these are self-checks of an initial review, not independent acceptance.\par
\end{document}
